% =====================================================================
% Segunda parte – 7. Conclusiones de las pruebas
% =====================================================================
\section{Conclusiones de las pruebas}

\paragraph{Unitarias.} Las 58 pruebas unitarias pasan. La caja blanca cubre todos los caminos factibles de los tres métodos con decisiones ($V(G) = 8$, 5 y 5) y la caja negra cubre las clases de equivalencia y los valores límite de los tres esquemas. Fueron estas pruebas las que encontraron la cuota negativa (D-01) y las validaciones ausentes (D-02 y D-03).

\paragraph{Integración.} La integración incremental ascendente, con pytest y Cypress, pasó de 6 de 11 casos aprobados a 18 de 18. Al ampliar las pruebas aparecieron dos defectos nuevos (D-11 y D-12), que se corrigieron con su prueba de regresión. La cobertura del backend es del 95\,\%.

\paragraph{Sistema.} Se eliminaron las credenciales codificadas (D-06), las entradas inválidas de los formularios ya no producen errores 500 (ZAP aún obtuvo uno con contraseñas de más de 72 bytes, límite de bcrypt) y la carga de 1\,500 peticiones con 50 usuarios terminó sin errores y con un percentil 90 de 22 ms en local. El riesgo principal sigue siendo la disponibilidad de la API en el plan gratuito de Render (D-09): la mitigación está lista, pero falta medirla en producción, igual que repetir JMeter sobre la API desplegada. ZAP y SonarQube se repitieron el 26/11/2026.

\paragraph{Aceptación.} Los dos casos que fallaban (PA-06 y PA-07) pasan después de las correcciones. El guion cubre ahora los 20 RF con 17 casos. En la sesión con el usuario final sobre la aplicación desplegada se aceptaron los 17 casos y los 20 RF (100\,\%), con tres observaciones que quedan como mejoras futuras o ya están mitigadas.
